\subsection{Background}

Recent advances in Artificial Intelligence (AI), particularly in Text-to-Speech (TTS) and Voice Conversion (VC) technologies, have significantly improved the quality and realism of synthetic speech generation. The adoption of deep neural networks, self-supervised speech representations, and neural vocoders has enabled the synthesis of speech that closely resembles natural human speech in terms of intelligibility, speaker characteristics, and prosodic features.

Research in neural speech synthesis has evolved rapidly over the last decade. Early end-to-end architectures such as Tacotron \cite{tacotron2017} demonstrated the feasibility of generating natural-sounding speech directly from text. Subsequent developments introduced improvements in prosody modeling \cite{prosody2018}, style control \cite{styletokens2018}, multi-speaker synthesis and voice transfer \cite{transfer2018}, further enhancing the realism and flexibility of generated speech. More recently, voice cloning systems have shown the ability to reproduce a target speaker’s voice using only a limited amount of reference audio.

As synthetic speech quality continues to improve, distinguishing between human and AI-generated speech has become increasingly challenging. This has led to growing research interest in Audio Deepfake Detection (ADD), which focuses on identifying artifacts and inconsistencies introduced during the speech generation process. Existing ADD approaches utilize various signal processing, machine learning, and deep learning techniques to differentiate synthetic speech from authentic human speech. Public benchmark datasets and evaluation challenges have further accelerated research in this area, resulting in substantial improvements in detection performance under controlled experimental conditions.
